\section{SPICE}\label{spice}\index{spice@SPICE}

\textbf{Keywords:} SPICE, Sources, Passives, Transistor Models, BSIM,
Foundries, Unit Transistors, gm/ID

\section{Simulation Program with Integrated Circuit
Emphasis}\label{simulation-program-with-integrated-circuit-emphasis}

To manufacture an integrated circuit we have to be able to predict how
it's going to work. The only way to predict is to rely on our knowledge
of physics, and build models of the real world in our computers.

One simulation strategy for a model of the real world, which absolutely
every single integrated circuit in the world has used to come into
existence, is SPICE.

Published in 1973 by Nagel and Pederson

\href{https://www2.eecs.berkeley.edu/Pubs/TechRpts/1973/ERL-m-382.pdf}{SPICE
(Simulation Program with Integrated Circuit Emphasis)}


{
\centering
\aicfig{media/nagel.png}
}

\emph{Figure 1: Title page of Nagel and Pederson's 1973 SPICE paper from
the 16th Midwest Symposium on Circuit Theory. Source: L. W. Nagel,
SPICE2 memorandum ERL-M520, UC Berkeley, 1975}

\subsection{Today}\label{today}\index{today@Today}

There are multiple SPICE programs that has been written, but they all
work in a similar fashion. There are expensive ones, closed source, and
open source.

Some are better at dealing with complex circuits, some are faster, and
some are more accurate. If you don't have money, then start with
ngspice.

\textbf{Commercial}
\href{https://www.cadence.com/ko_KR/home/tools/custom-ic-analog-rf-design/circuit-simulation/spectre-simulation-platform.html}{Cadence
Spectre} \href{https://eda.sw.siemens.com/en-US/ic/eldo/}{Siemens Eldo}
\href{https://www.synopsys.com/implementation-and-signoff/ams-simulation/primesim-hspice.html}{Synopsys
HSPICE}

\textbf{Free} \href{http://aimspice.com}{Aimspice}
\href{https://www.analog.com/en/design-center/design-tools-and-calculators/ltspice-simulator.html}{Analog
Devices LTspice}

\textbf{Open Source} \href{http://ngspice.sourceforge.net}{ngspice}

\subsection{But}\label{but}

All SPICE simulators understand the same language (yes, even spectre can
speak SPICE). We write our testbenches in a text file, and give it to
the SPICE program. That's the same for all programs. Some may have built
fancy GUI's to hide the fact that we're really writing text files, but
text files is what is under the hood.

Pretty much the same usage model as 50-odd years ago

\begin{Shaded}
\begin{Highlighting}[]
\OperatorTok{\textless{}}\NormalTok{spice }\ExtensionTok{program}\OperatorTok{\textgreater{}}\NormalTok{ testbench.cir}
\end{Highlighting}
\end{Shaded}

for example

\begin{Shaded}
\begin{Highlighting}[]
\ExtensionTok{ngspice}\NormalTok{ testbench.cir}
\end{Highlighting}
\end{Shaded}

Or in the most expensive analog tool (Cadence Spectre)

\begin{Shaded}
\begin{Highlighting}[]
\ExtensionTok{spectre}\NormalTok{  input.scs  +escchars +log ../psf/spectre.out }
\ExtensionTok{{-}format}\NormalTok{ psfxl }\AttributeTok{{-}raw}\NormalTok{ ../psf   +aps +lqtimeout 900 }\AttributeTok{{-}maxw}\NormalTok{ 5 }
\ExtensionTok{{-}maxn}\NormalTok{ 5 }\AttributeTok{{-}env}\NormalTok{ ade  }\AttributeTok{{-}ahdllibdir} 
\ExtensionTok{/tmp/wulff/virtuoso/TB\_SUN\_BIAS\_GF130N/TB\_SUN\_BIAS/maestro/}
\ExtensionTok{results/maestro/Interactive.15/sharedData/CDS/ahdl/input.ahdlSimDB} 
\ExtensionTok{+logstatus} 
\end{Highlighting}
\end{Shaded}

The expensive tools have built graphical user interface around the SPICE
simulator to make it easier to run multiple scenarios.

\begin{longtable}[]{@{}lllll@{}}
\toprule\noalign{}
Corner & Typical & Fast & Slow & All \\
\midrule\noalign{}
\endhead
\bottomrule\noalign{}
\endlastfoot
Mosfet & Mtt & Mff & Mss & Mff,Mfs,Msf,Mss \\
Resistor & Rt & Rl & Rh & Rl,Rh \\
Capacitors & Ct & Cl & Ch & Cl,Ch \\
Diode & Dt & Df & Ds & Df,Ds \\
Bipolar & Bt & Bf & Bs & Bf,Bs \\
Temperature & Tt & Th,Tl & Th,Tl & Th,Tl \\
Voltage & Vt & Vh,Vl & Vh,Vl & Vh,Vl \\
\end{longtable}


{
\centering
\aicfig{media/assembler.png}
}

\emph{Figure 2: Cadence Virtuoso ADE Assembler, showing the corner
definitions on the left and a pass/fail table of simulated
specifications on the right}

I'm a fan of launching multiple simulations from the command line. I
don't like GUI's. As such, I wrote
\href{https://github.com/wulffern/cicsim/tree/main}{cicsim}, and that's
what I use in the video and demo.

\subsection{Sources}\label{sources}\index{sources@Sources}

The SPICE language is a set of conventions for how to write the text
files. In general, it's one line, one command (although, lines can be
continued with a +).

I'm not going to go through an extensive tutorial in this document, and
there are dialects with different SPICE programs. You'll find more info
at \href{http://ngspice.sourceforge.net/docs.html}{ngspice}

\subsubsection{Independent current
sources}\label{independent-current-sources}

\emph{Infinite output impedance, changing voltage does not change
current}

\begin{Shaded}
\begin{Highlighting}[]
\ExtensionTok{I}\OperatorTok{\textless{}}\NormalTok{name}\OperatorTok{\textgreater{}} \OperatorTok{\textless{}}\NormalTok{from}\OperatorTok{\textgreater{}} \OperatorTok{\textless{}}\NormalTok{to}\OperatorTok{\textgreater{}}\NormalTok{ dc }\OperatorTok{\textless{}}\NormalTok{number}\OperatorTok{\textgreater{}}\NormalTok{ ac }\OperatorTok{\textless{}}\NormalTok{number}\OperatorTok{\textgreater{}}

\ExtensionTok{I1}\NormalTok{ 0 VDN dc In}
\ExtensionTok{I2}\NormalTok{ VDP 0 dc Ip}
\end{Highlighting}
\end{Shaded}

\subsubsection{Independent voltage
source}\label{independent-voltage-source}

\emph{Zero output impedance, changing current does not change voltage}

\begin{Shaded}
\begin{Highlighting}[]
\ExtensionTok{V}\OperatorTok{\textless{}}\NormalTok{name}\OperatorTok{\textgreater{}} \OperatorTok{\textless{}}\NormalTok{+}\OperatorTok{\textgreater{}} \OperatorTok{\textless{}}\NormalTok{{-}}\OperatorTok{\textgreater{}}\NormalTok{ dc }\OperatorTok{\textless{}}\NormalTok{number}\OperatorTok{\textgreater{}}\NormalTok{ ac }\OperatorTok{\textless{}}\NormalTok{number}\OperatorTok{\textgreater{}}

\ExtensionTok{V2}\NormalTok{ VSS 0 dc 0}
\ExtensionTok{V1}\NormalTok{ VDD 0 dc 1.5}
\end{Highlighting}
\end{Shaded}

\subsection{Passives}\label{passives}\index{passives@Passives}

Resistors

\begin{Shaded}
\begin{Highlighting}[]
\ExtensionTok{R}\OperatorTok{\textless{}}\NormalTok{name}\OperatorTok{\textgreater{}} \OperatorTok{\textless{}}\NormalTok{node }\DecValTok{1}\OperatorTok{\textgreater{}} \OperatorTok{\textless{}}\NormalTok{node }\DecValTok{2}\OperatorTok{\textgreater{}} \OperatorTok{\textless{}}\NormalTok{value}\OperatorTok{\textgreater{}}

\ExtensionTok{R1}\NormalTok{ N1 N2 10k}
\ExtensionTok{R2}\NormalTok{ N2 N3 1Meg}
\ExtensionTok{R3}\NormalTok{ N3 N4 1G}
\ExtensionTok{R4}\NormalTok{ N4 N5 1T}
\end{Highlighting}
\end{Shaded}

Capacitors

\begin{Shaded}
\begin{Highlighting}[]
\ExtensionTok{C}\OperatorTok{\textless{}}\NormalTok{name}\OperatorTok{\textgreater{}} \OperatorTok{\textless{}}\NormalTok{node }\DecValTok{1}\OperatorTok{\textgreater{}} \OperatorTok{\textless{}}\NormalTok{node }\DecValTok{2}\OperatorTok{\textgreater{}} \OperatorTok{\textless{}}\NormalTok{value}\OperatorTok{\textgreater{}}

\ExtensionTok{C1}\NormalTok{ N1 N2 1a}
\ExtensionTok{C2}\NormalTok{ N1 N2 1f}
\ExtensionTok{C4}\NormalTok{ N1 N2 1p}
\ExtensionTok{C3}\NormalTok{ N1 N2 1n}
\ExtensionTok{C5}\NormalTok{ N1 N2 1u}
\end{Highlighting}
\end{Shaded}

\subsection{Transistor Models}\label{transistor-models}\index{transistor models@Transistor models}

Needs a model file describing the transistor model

BSIM (Berkeley Short-channel IGFET Model)
\url{http://bsim.berkeley.edu/models/bsim4/}


{
\centering
\aicfig{media/fig_transistor.png}
}

\emph{Figure 3: Circuit symbol for an NMOS transistor with its gate,
drain and source terminals}

284 parameters in
\href{http://www-device.eecs.berkeley.edu/~bsim/Files/BSIM4/BSIM480/BSIM480_Manual.pdf}{BSIM
4.5}

\begin{Shaded}
\begin{Highlighting}[]
\AttributeTok{.}\ConstantTok{MODEL} \ConstantTok{N1} \ConstantTok{NMOS} \ConstantTok{LEVEL}\OperatorTok{=}\DecValTok{14} \ConstantTok{VERSION}\OperatorTok{=}\FloatTok{4.5}\AttributeTok{.}\DecValTok{0} \ConstantTok{BINUNIT}\OperatorTok{=}\DecValTok{1} 
\ConstantTok{PARAMCHK}\OperatorTok{=}\DecValTok{1} \ConstantTok{MOBMOD}\OperatorTok{=}\DecValTok{0} \ConstantTok{CAPMOD}\OperatorTok{=}\DecValTok{2} \ConstantTok{IGCMOD}\OperatorTok{=}\DecValTok{1} \ConstantTok{IGBMOD}\OperatorTok{=}\DecValTok{1} 
\ConstantTok{GEOMOD}\OperatorTok{=}\DecValTok{1}  \ConstantTok{DIOMOD}\OperatorTok{=}\DecValTok{1} \ConstantTok{RDSMOD}\OperatorTok{=}\DecValTok{0} \ConstantTok{RBODYMOD}\OperatorTok{=}\DecValTok{0} \ConstantTok{RGATEMOD}\OperatorTok{=}\DecValTok{3}
\ConstantTok{PERMOD}\OperatorTok{=}\DecValTok{1} \ConstantTok{ACNQSMOD}\OperatorTok{=}\DecValTok{0} \ConstantTok{TRNQSMOD}\OperatorTok{=}\DecValTok{0} \ConstantTok{TEMPMOD}\OperatorTok{=}\DecValTok{0}  \ConstantTok{TNOM}\OperatorTok{=}\DecValTok{27} 
\ConstantTok{TOXE}\OperatorTok{=}\FloatTok{1.8E{-}009} \ConstantTok{TOXP}\OperatorTok{=}\FloatTok{10E{-}010} \ConstantTok{TOXM}\OperatorTok{=}\FloatTok{1.8E{-}009}  \ConstantTok{DTOX}\OperatorTok{=}\FloatTok{8E{-}10} 
\ConstantTok{EPSROX}\OperatorTok{=}\FloatTok{3.9} \ConstantTok{WINT}\OperatorTok{=}\FloatTok{5E{-}009} \ConstantTok{LINT}\OperatorTok{=}\FloatTok{1E{-}009} \ConstantTok{LL}\OperatorTok{=}\DecValTok{0} \ConstantTok{WL}\OperatorTok{=}\DecValTok{0} \ConstantTok{LLN}\OperatorTok{=}\DecValTok{1} 
\ConstantTok{WLN}\OperatorTok{=}\DecValTok{1}  \ConstantTok{LW}\OperatorTok{=}\DecValTok{0} \ConstantTok{WW}\OperatorTok{=}\DecValTok{0} \ConstantTok{LWN}\OperatorTok{=}\DecValTok{1} \ConstantTok{WWN}\OperatorTok{=}\DecValTok{1}  \ConstantTok{LWL}\OperatorTok{=}\DecValTok{0} \ConstantTok{WWL}\OperatorTok{=}\DecValTok{0} \ConstantTok{XPART}\OperatorTok{=}\DecValTok{0}
\ConstantTok{TOXREF}\OperatorTok{=}\FloatTok{1.4E{-}009}  \ConstantTok{SAREF}\OperatorTok{=}\FloatTok{5E{-}6} \ConstantTok{SBREF}\OperatorTok{=}\FloatTok{5E{-}6} \ConstantTok{WLOD}\OperatorTok{=}\FloatTok{2E{-}6} 
\ConstantTok{KU0}\OperatorTok{={-}}\FloatTok{4E{-}6}  \ConstantTok{KVSAT}\OperatorTok{=}\FloatTok{0.2} \ConstantTok{KVTH0}\OperatorTok{={-}}\FloatTok{2E{-}8} \ConstantTok{TKU0}\OperatorTok{=}\FloatTok{0.0} \ConstantTok{LLODKU0}\OperatorTok{=}\FloatTok{1.1}  
\ConstantTok{WLODKU0}\OperatorTok{=}\FloatTok{1.1} \ConstantTok{LLODVTH}\OperatorTok{=}\FloatTok{1.0} \ConstantTok{WLODVTH}\OperatorTok{=}\FloatTok{1.0} \ConstantTok{LKU0}\OperatorTok{=}\FloatTok{1E{-}6}  
\ConstantTok{WKU0}\OperatorTok{=}\FloatTok{1E{-}6} \ConstantTok{PKU0}\OperatorTok{=}\FloatTok{0.0} \ConstantTok{LKVTH0}\OperatorTok{=}\FloatTok{1.1E{-}6} \ConstantTok{WKVTH0}\OperatorTok{=}\FloatTok{1.1E{-}6}  
\ConstantTok{PKVTH0}\OperatorTok{=}\FloatTok{0.0} \ConstantTok{STK2}\OperatorTok{=}\FloatTok{0.0} \ConstantTok{LODK2}\OperatorTok{=}\FloatTok{1.0} \ConstantTok{STETA0}\OperatorTok{=}\FloatTok{0.0}  \ConstantTok{LODETA0}\OperatorTok{=}\FloatTok{1.0}  
\ConstantTok{LAMBDA}\OperatorTok{=}\FloatTok{4E{-}10}  \ConstantTok{VSAT}\OperatorTok{=}\FloatTok{1.1}\ErrorTok{E} \BaseNTok{005} \ConstantTok{VTL}\OperatorTok{=}\FloatTok{2.0E5} \ConstantTok{XN}\OperatorTok{=}\FloatTok{6.0} \ConstantTok{LC}\OperatorTok{=}\FloatTok{5E{-}9} 
\ConstantTok{RNOIA}\OperatorTok{=}\FloatTok{0.577} \ConstantTok{RNOIB}\OperatorTok{=}\FloatTok{0.37} \ConstantTok{LINTNOI}\OperatorTok{=}\FloatTok{1E{-}009}  \ConstantTok{WPEMOD}\OperatorTok{=}\DecValTok{0} 
\ConstantTok{WEB}\OperatorTok{=}\FloatTok{0.0} \ConstantTok{WEC}\OperatorTok{=}\FloatTok{0.0} \ConstantTok{KVTH0WE}\OperatorTok{=}\FloatTok{1.0}  \ConstantTok{K2WE}\OperatorTok{=}\FloatTok{1.0} \ConstantTok{KU0WE}\OperatorTok{=}\FloatTok{1.0} \ConstantTok{SCREF}\OperatorTok{=}\FloatTok{5.0E{-}6}
\ConstantTok{TVOFF}\OperatorTok{=}\FloatTok{0.0} \ConstantTok{TVFBSDOFF}\OperatorTok{=}\FloatTok{0.0}  \ConstantTok{VTH0}\OperatorTok{=}\FloatTok{0.25}  \ConstantTok{K1}\OperatorTok{=}\FloatTok{0.35} \ConstantTok{K2}\OperatorTok{=}\FloatTok{0.05}
\ConstantTok{K3}\OperatorTok{=}\DecValTok{0}  \ConstantTok{K3B}\OperatorTok{=}\DecValTok{0} \ConstantTok{W0}\OperatorTok{=}\FloatTok{2.5E{-}006} \ConstantTok{DVT0}\OperatorTok{=}\FloatTok{1.8} \ConstantTok{DVT1}\OperatorTok{=}\FloatTok{0.52}  \ConstantTok{DVT2}\OperatorTok{={-}}\FloatTok{0.032} 
\ConstantTok{DVT0W}\OperatorTok{=}\DecValTok{0} \ConstantTok{DVT1W}\OperatorTok{=}\DecValTok{0} \ConstantTok{DVT2W}\OperatorTok{=}\DecValTok{0}  \ConstantTok{DSUB}\OperatorTok{=}\DecValTok{2} \ConstantTok{MINV}\OperatorTok{=}\FloatTok{0.05} \ConstantTok{VOFFL}\OperatorTok{=}\DecValTok{0} 
\ConstantTok{DVTP0}\OperatorTok{=}\FloatTok{1E{-}007}  \ConstantTok{DVTP1}\OperatorTok{=}\FloatTok{0.05} \ConstantTok{LPE0}\OperatorTok{=}\FloatTok{5.75E{-}008} \ConstantTok{LPEB}\OperatorTok{=}\FloatTok{2.3E{-}010} 
\ConstantTok{XJ}\OperatorTok{=}\FloatTok{2E{-}008}  \ConstantTok{NGATE}\OperatorTok{=}\ErrorTok{5E} \BaseNTok{020} \ConstantTok{NDEP}\OperatorTok{=}\FloatTok{2.8}\ErrorTok{E} \BaseNTok{01}\ErrorTok{8} \ConstantTok{NSD}\OperatorTok{=}\ErrorTok{1E} \BaseNTok{020} \ConstantTok{PHIN}\OperatorTok{=}\DecValTok{0} 
\ConstantTok{CDSC}\OperatorTok{=}\FloatTok{0.0002} \ConstantTok{CDSCB}\OperatorTok{=}\DecValTok{0} \ConstantTok{CDSCD}\OperatorTok{=}\DecValTok{0} \ConstantTok{CIT}\OperatorTok{=}\DecValTok{0}  \ConstantTok{VOFF}\OperatorTok{={-}}\FloatTok{0.15} \ConstantTok{NFACTOR}\OperatorTok{=}\FloatTok{1.2} 
\ConstantTok{ETA0}\OperatorTok{=}\FloatTok{0.05} \ConstantTok{ETAB}\OperatorTok{=}\DecValTok{0}  \ConstantTok{UC}\OperatorTok{={-}}\FloatTok{3E{-}011}  \ConstantTok{VFB}\OperatorTok{={-}}\FloatTok{0.55} \ConstantTok{U0}\OperatorTok{=}\FloatTok{0.032} 
\ConstantTok{UA}\OperatorTok{=}\FloatTok{5.0E{-}011} \ConstantTok{UB}\OperatorTok{=}\FloatTok{3.5E{-}018}  \ConstantTok{A0}\OperatorTok{=}\DecValTok{2} \ConstantTok{AGS}\OperatorTok{=}\FloatTok{1E{-}020} \ConstantTok{A1}\OperatorTok{=}\DecValTok{0} \ConstantTok{A2}\OperatorTok{=}\DecValTok{1} 
\ConstantTok{B0}\OperatorTok{={-}}\FloatTok{1E{-}020} \ConstantTok{B1}\OperatorTok{=}\DecValTok{0}  \ConstantTok{KETA}\OperatorTok{=}\FloatTok{0.04} \ConstantTok{DWG}\OperatorTok{=}\DecValTok{0} \ConstantTok{DWB}\OperatorTok{=}\DecValTok{0} \ConstantTok{PCLM}\OperatorTok{=}\FloatTok{0.08}
\ConstantTok{PDIBLC1}\OperatorTok{=}\FloatTok{0.028} \ConstantTok{PDIBLC2}\OperatorTok{=}\FloatTok{0.022} \ConstantTok{PDIBLCB}\OperatorTok{={-}}\FloatTok{0.005} \ConstantTok{DROUT}\OperatorTok{=}\FloatTok{0.45}  
\ConstantTok{PVAG}\OperatorTok{=}\FloatTok{1E{-}020} \ConstantTok{DELTA}\OperatorTok{=}\FloatTok{0.01} \ConstantTok{PSCBE1}\OperatorTok{=}\FloatTok{8.14}\ErrorTok{E} \BaseNTok{00}\ErrorTok{8} \ConstantTok{PSCBE2}\OperatorTok{=}\FloatTok{5E{-}008}  
\ConstantTok{RSH}\OperatorTok{=}\DecValTok{0} \ConstantTok{RDSW}\OperatorTok{=}\DecValTok{0} \ConstantTok{RSW}\OperatorTok{=}\DecValTok{0} \ConstantTok{RDW}\OperatorTok{=}\DecValTok{0} \ConstantTok{FPROUT}\OperatorTok{=}\FloatTok{0.2} \ConstantTok{PDITS}\OperatorTok{=}\FloatTok{0.2} \ConstantTok{PDITSD}\OperatorTok{=}\FloatTok{0.23} 
\ConstantTok{PDITSL}\OperatorTok{=}\FloatTok{2.3}\ErrorTok{E} \BaseNTok{006}  \ConstantTok{RSH}\OperatorTok{=}\DecValTok{0} \ConstantTok{RDSW}\OperatorTok{=}\DecValTok{50} \ConstantTok{RSW}\OperatorTok{=}\DecValTok{150}
\ConstantTok{RDW}\OperatorTok{=}\DecValTok{150}  \ConstantTok{RDSWMIN}\OperatorTok{=}\DecValTok{0} \ConstantTok{RDWMIN}\OperatorTok{=}\DecValTok{0} \ConstantTok{RSWMIN}\OperatorTok{=}\DecValTok{0} \ConstantTok{PRWG}\OperatorTok{=}\DecValTok{0}  \ConstantTok{PRWB}\OperatorTok{=}\FloatTok{6.8E{-}011} 
\ConstantTok{WR}\OperatorTok{=}\DecValTok{1} \ConstantTok{ALPHA0}\OperatorTok{=}\FloatTok{0.074} \ConstantTok{ALPHA1}\OperatorTok{=}\FloatTok{0.005}  \ConstantTok{BETA0}\OperatorTok{=}\DecValTok{30} \ConstantTok{AGIDL}\OperatorTok{=}\FloatTok{0.0002} 
\ConstantTok{BGIDL}\OperatorTok{=}\FloatTok{2.1}\ErrorTok{E} \BaseNTok{00}\ErrorTok{9} \ConstantTok{CGIDL}\OperatorTok{=}\FloatTok{0.0002} \ConstantTok{EGIDL}\OperatorTok{=}\FloatTok{0.8}  \ConstantTok{AIGBACC}\OperatorTok{=}\FloatTok{0.012} 
\ConstantTok{BIGBACC}\OperatorTok{=}\FloatTok{0.0028} \ConstantTok{CIGBACC}\OperatorTok{=}\FloatTok{0.002}  \ConstantTok{NIGBACC}\OperatorTok{=}\DecValTok{1} \ConstantTok{AIGBINV}\OperatorTok{=}\FloatTok{0.014} 
\ConstantTok{BIGBINV}\OperatorTok{=}\FloatTok{0.004} \ConstantTok{CIGBINV}\OperatorTok{=}\FloatTok{0.004}  \ConstantTok{EIGBINV}\OperatorTok{=}\FloatTok{1.1} \ConstantTok{NIGBINV}\OperatorTok{=}\DecValTok{3} \ConstantTok{AIGC}\OperatorTok{=}\FloatTok{0.012}
\ConstantTok{BIGC}\OperatorTok{=}\FloatTok{0.0028}  \ConstantTok{CIGC}\OperatorTok{=}\FloatTok{0.002} \ConstantTok{AIGSD}\OperatorTok{=}\FloatTok{0.012} \ConstantTok{BIGSD}\OperatorTok{=}\FloatTok{0.0028} \ConstantTok{CIGSD}\OperatorTok{=}\FloatTok{0.002}  \ConstantTok{NIGC}\OperatorTok{=}\DecValTok{1}
\ConstantTok{POXEDGE}\OperatorTok{=}\DecValTok{1} \ConstantTok{PIGCD}\OperatorTok{=}\DecValTok{1} \ConstantTok{NTOX}\OperatorTok{=}\DecValTok{1}  \ConstantTok{VFBSDOFF}\OperatorTok{=}\FloatTok{0.0}  \ConstantTok{XRCRG1}\OperatorTok{=}\DecValTok{12} \ConstantTok{XRCRG2}\OperatorTok{=}\DecValTok{5}  
\ConstantTok{CGSO}\OperatorTok{=}\FloatTok{6.238E{-}010} \ConstantTok{CGDO}\OperatorTok{=}\FloatTok{6.238E{-}010} \ConstantTok{CGBO}\OperatorTok{=}\FloatTok{2.56E{-}011} \ConstantTok{CGDL}\OperatorTok{=}\FloatTok{2.495E{-}10}  
\ConstantTok{CGSL}\OperatorTok{=}\FloatTok{2.495E{-}10} \ConstantTok{CKAPPAS}\OperatorTok{=}\FloatTok{0.03} \ConstantTok{CKAPPAD}\OperatorTok{=}\FloatTok{0.03} \ConstantTok{ACDE}\OperatorTok{=}\DecValTok{1}  \ConstantTok{MOIN}\OperatorTok{=}\DecValTok{15} 
\ConstantTok{NOFF}\OperatorTok{=}\FloatTok{0.9} \ConstantTok{VOFFCV}\OperatorTok{=}\FloatTok{0.02}  \ConstantTok{KT1}\OperatorTok{={-}}\FloatTok{0.37} \ConstantTok{KT1L}\OperatorTok{=}\FloatTok{0.0} \ConstantTok{KT2}\OperatorTok{={-}}\FloatTok{0.042} \ConstantTok{UTE}\OperatorTok{={-}}\FloatTok{1.5}  
\ConstantTok{UA1}\OperatorTok{=}\FloatTok{1E{-}009} \ConstantTok{UB1}\OperatorTok{={-}}\FloatTok{3.5E{-}019} \ConstantTok{UC1}\OperatorTok{=}\DecValTok{0} \ConstantTok{PRT}\OperatorTok{=}\DecValTok{0} \ConstantTok{AT}\OperatorTok{=}\DecValTok{53000}  \ConstantTok{FNOIMOD}\OperatorTok{=}\DecValTok{1} 
\ConstantTok{TNOIMOD}\OperatorTok{=}\DecValTok{0}  \ConstantTok{JSS}\OperatorTok{=}\FloatTok{0.0001} \ConstantTok{JSWS}\OperatorTok{=}\FloatTok{1E{-}011} \ConstantTok{JSWGS}\OperatorTok{=}\FloatTok{1E{-}010} \ConstantTok{NJS}\OperatorTok{=}\DecValTok{1}
\ConstantTok{IJTHSFWD}\OperatorTok{=}\FloatTok{0.01} \ConstantTok{IJTHSREV}\OperatorTok{=}\FloatTok{0.001} \ConstantTok{BVS}\OperatorTok{=}\DecValTok{10} \ConstantTok{XJBVS}\OperatorTok{=}\DecValTok{1}  \ConstantTok{JSD}\OperatorTok{=}\FloatTok{0.0001} 
\ConstantTok{JSWD}\OperatorTok{=}\FloatTok{1E{-}011} \ConstantTok{JSWGD}\OperatorTok{=}\FloatTok{1E{-}010} \ConstantTok{NJD}\OperatorTok{=}\DecValTok{1}  \ConstantTok{IJTHDFWD}\OperatorTok{=}\FloatTok{0.01} \ConstantTok{IJTHDREV}\OperatorTok{=}\FloatTok{0.001} 
\ConstantTok{BVD}\OperatorTok{=}\DecValTok{10} \ConstantTok{XJBVD}\OperatorTok{=}\DecValTok{1}  \ConstantTok{PBS}\OperatorTok{=}\DecValTok{1} \ConstantTok{CJS}\OperatorTok{=}\FloatTok{0.0005} \ConstantTok{MJS}\OperatorTok{=}\FloatTok{0.5} \ConstantTok{PBSWS}\OperatorTok{=}\DecValTok{1}  \ConstantTok{CJSWS}\OperatorTok{=}\FloatTok{5E{-}010} 
\ConstantTok{MJSWS}\OperatorTok{=}\FloatTok{0.33} \ConstantTok{PBSWGS}\OperatorTok{=}\DecValTok{1} \ConstantTok{CJSWGS}\OperatorTok{=}\FloatTok{3E{-}010}  \ConstantTok{MJSWGS}\OperatorTok{=}\FloatTok{0.33} \ConstantTok{PBD}\OperatorTok{=}\DecValTok{1} \ConstantTok{CJD}\OperatorTok{=}\FloatTok{0.0005} 
\ConstantTok{MJD}\OperatorTok{=}\FloatTok{0.5}  \ConstantTok{PBSWD}\OperatorTok{=}\DecValTok{1} \ConstantTok{CJSWD}\OperatorTok{=}\FloatTok{5E{-}010} \ConstantTok{MJSWD}\OperatorTok{=}\FloatTok{0.33} \ConstantTok{PBSWGD}\OperatorTok{=}\DecValTok{1}
\ConstantTok{CJSWGD}\OperatorTok{=}\FloatTok{5E{-}010}\ErrorTok{MJSWGD}\OperatorTok{=}\FloatTok{0.33} \ConstantTok{TPB}\OperatorTok{=}\FloatTok{0.005} \ConstantTok{TCJ}\OperatorTok{=}\FloatTok{0.001} \ConstantTok{TPBSW}\OperatorTok{=}\FloatTok{0.005} 
\ConstantTok{TCJSW}\OperatorTok{=}\FloatTok{0.001} \ConstantTok{TPBSWG}\OperatorTok{=}\FloatTok{0.005} \ConstantTok{TCJSWG}\OperatorTok{=}\FloatTok{0.001}  \ConstantTok{XTIS}\OperatorTok{=}\DecValTok{3} \ConstantTok{XTID}\OperatorTok{=}\DecValTok{3}  \ConstantTok{DMCG}\OperatorTok{=}\FloatTok{0E{-}006} 
\ConstantTok{DMCI}\OperatorTok{=}\FloatTok{0E{-}006} \ConstantTok{DMDG}\OperatorTok{=}\FloatTok{0E{-}006} \ConstantTok{DMCGT}\OperatorTok{=}\FloatTok{0E{-}007}  \ConstantTok{DWJ}\OperatorTok{=}\FloatTok{0.0E{-}008} \ConstantTok{XGW}\OperatorTok{=}\FloatTok{0E{-}007}
\ConstantTok{XGL}\OperatorTok{=}\FloatTok{0E{-}008}  \ConstantTok{RSHG}\OperatorTok{=}\FloatTok{0.4} \ConstantTok{GBMIN}\OperatorTok{=}\FloatTok{1E{-}010} \ConstantTok{RBPB}\OperatorTok{=}\DecValTok{5} \ConstantTok{RBPD}\OperatorTok{=}\DecValTok{15}  \ConstantTok{RBPS}\OperatorTok{=}\DecValTok{15} \ConstantTok{RBDB}\OperatorTok{=}\DecValTok{15} 
\ConstantTok{RBSB}\OperatorTok{=}\DecValTok{15} \ConstantTok{NGCON}\OperatorTok{=}\DecValTok{1} \ConstantTok{JTSS}\OperatorTok{=}\FloatTok{1E{-}4} \ConstantTok{JTSD}\OperatorTok{=}\FloatTok{1E{-}4} \ConstantTok{JTSSWS}\OperatorTok{=}\FloatTok{1E{-}10} \ConstantTok{JTSSWD}\OperatorTok{=}\FloatTok{1E{-}10} 
\ConstantTok{JTSSWGS}\OperatorTok{=}\FloatTok{1E{-}7} \ConstantTok{JTSSWGD}\OperatorTok{=}\FloatTok{1E{-}7}  \ConstantTok{NJTS}\OperatorTok{=}\FloatTok{20.0} \ConstantTok{NJTSSW}\OperatorTok{=}\DecValTok{20} \ConstantTok{NJTSSWG}\OperatorTok{=}\DecValTok{6} 
\ConstantTok{VTSS}\OperatorTok{=}\DecValTok{10} \ConstantTok{VTSD}\OperatorTok{=}\DecValTok{10} \ConstantTok{VTSSWS}\OperatorTok{=}\DecValTok{10} \ConstantTok{VTSSWD}\OperatorTok{=}\DecValTok{10}  \ConstantTok{VTSSWGS}\OperatorTok{=}\DecValTok{2} \ConstantTok{VTSSWGD}\OperatorTok{=}\DecValTok{2}
\ConstantTok{XTSS}\OperatorTok{=}\FloatTok{0.02} \ConstantTok{XTSD}\OperatorTok{=}\FloatTok{0.02} \ConstantTok{XTSSWS}\OperatorTok{=}\FloatTok{0.02} \ConstantTok{XTSSWD}\OperatorTok{=}\FloatTok{0.02} \ConstantTok{XTSSWGS}\OperatorTok{=}\FloatTok{0.02} 
\ConstantTok{XTSSWGD}\OperatorTok{=}\FloatTok{0.02}
\end{Highlighting}
\end{Shaded}

\subsection{Transistors}\label{transistors}\index{transistors@Transistors}

\begin{Shaded}
\begin{Highlighting}[]

\ExtensionTok{M}\OperatorTok{\textless{}}\NormalTok{name}\OperatorTok{\textgreater{}} \OperatorTok{\textless{}}\NormalTok{drain}\OperatorTok{\textgreater{}} \OperatorTok{\textless{}}\NormalTok{gate}\OperatorTok{\textgreater{}} \OperatorTok{\textless{}}\NormalTok{source}\OperatorTok{\textgreater{}} \OperatorTok{\textless{}}\NormalTok{bulk}\OperatorTok{\textgreater{}} \OperatorTok{\textless{}}\NormalTok{modelname}\OperatorTok{\textgreater{}} \PreprocessorTok{[}\SpecialStringTok{parameters}\PreprocessorTok{]}


\ExtensionTok{M1}\NormalTok{ VDN VDN VSS VSS nmos W=0.6u L=0.15u}
\ExtensionTok{M2}\NormalTok{ VDP VDP VDD VDD pmos W=0.6u L=0.15u}

\end{Highlighting}
\end{Shaded}

\subsection{Foundries}\label{foundries}\index{foundries@Foundries}

Each foundry has their own SPICE models bacause the transistor
parameters depend on the exact physics of the technology!

\url{https://skywater-pdk.readthedocs.io/en/main/}

\section{Find right transistor sizes}\label{find-right-transistor-sizes}\index{find right transistor sizes@Find right transistor sizes}

Assume active (\[V_{ds} > V_{eff}\] in strong inversion, or
\[V_{ds} > 3 V_T\] in weak inversion). For diode connected transistors,
that is always true.

Weak inversion: \[ I_{D} = I_{D0} \frac{W}{L} e^{V_eff / n V_T} \],
\[V_{eff} \propto \ln{I_D} \]

Strong inversion:
\[ I_{D} = \frac{1}{2} \mu_n C_{ox} \frac{W}{L} V_{eff}^2\],
\[V_{eff} \propto \sqrt{I_D} \]

\textbf{Operating region for a diode connected transistor only depends
on the current}


{
\centering
\aicfig{media/trop.png}
}

\emph{Figure 4: Diode connected NMOS biased by a current source, where
the current alone sets the operating region}

\subsection{Use unit size transistors for analog
design}\label{use-unit-size-transistors-for-analog-design}

\[ W/L \approx \in[4, 6, 10] \], but should have space for two contacts

Use parallel transistors for larger W/L

Amplifiers \[\Rightarrow L \approx 1.2 \times L_{min} \]

Current mirrors \[\Rightarrow L \approx 4 \times L_{min} \]

Choose sizes that have been used by foundry for measurement to match
SPICE model

\subsection{What about gm/Id ?}\label{what-about-gmid}\index{what about gm/id ?@What about gm/Id ?}

Weak \[ \frac{g_m}{I_d} = \frac{1}{nV_T}\]

Strong \[ \frac{g_m}{I_d} = \frac{2}{V_{eff}}\]

\subsection{Characterize the
transistors}\label{characterize-the-transistors}

\url{http://analogicus.com/cnr_atr_sky130nm/mos/CNRATR_NCH_2C1F2.html}

\section{More information}\label{more-information}\index{more information@More information}

\href{http://ngspice.sourceforge.net/docs/ngspice-43-manual.pdf}{Ngspice
Manual}

\href{https://analogicus.com/aicex/started/}{Installing tools}

\section{Analog Design}\label{analog-design}\index{analog design@Analog design}

\begin{enumerate}
\def\labelenumi{\arabic{enumi}.}
\tightlist
\item
  \textbf{Define the problem, what are you trying to solve?}
\item
  \textbf{Find a circuit that can solve the problem (papers, books)}
\item
  \textbf{Find right transistor sizes. What transistors should be weak
  inversion, strong inversion, or don't care?}
\item
  \textbf{Check operating region of transistors (.op)}
\item
  \textbf{Check key parameters (.dc, .ac, .tran)}
\item
  \textbf{Check function. Exercise all inputs. Check all control
  signals}
\item
  Check key parameters in all corners. Check mismatch (Monte-Carlo
  simulation)
\item
  Do layout, and check it's error free. Run design rule checks (DRC).
  Check layout versus schematic (LVS)
\item
  Extract parasitics from layout. Resistance, capacitance, and
  inductance if necessary.
\item
  On extracted parasitic netlist, check key parameters in all corners
  and mismatch (if possible).
\item
  If everything works, then you're done.
\end{enumerate}

\emph{On failure, go back}

\section{Demo}\label{demo}

\url{https://github.com/analogicus/jnw_spice_sky130A/tree/main}

\section{Summary}\label{summary}

The one-page version of this chapter:

\begin{itemize}
\tightlist
\item
  SPICE has run the same way for fifty-odd years: a netlist in,
  operating points and waveforms out
\item
  The quartet to master: op, dc, ac and tran - everything else is
  decoration on those four
\item
  ngspice speaks the dialect this course uses, and the transistor models
  come from the PDK, not from the simulator
\end{itemize}

\section{Would you like to know
more?}\label{would-you-like-to-know-more}

Nagel's 1975 thesis, where SPICE comes from, and which still reads well

The ngspice manual, the reference for everything this chapter does
